\chapter{Ethische und sozialwissenschaftliche Perspektiven}

\section{Technikethik}

Beim Einsatz von technischen Systemen entstehen moralische Bedenken, welche ethische Vertretbarkeit von den technischen Systemen in Frage stellen. In Bezug auf \textit{KI} und \textit{DAS} entstehen drei wesentliche Ansätze.

(Hans Jonas - Prinzip Verantwortung) Langfristige Folgen beim Einsatz von diesen Systemen einzuschätzen und positiv entgegenzuwirken. In unserem Beispiel wird der ökologische Fußabdruck berücksichtigt.

(van de Poel - Werte-sensitive Technikgestaltung) Die Gestaltung von technischen Systemen soll mit Rücksicht auf gesellschaftliche Werte gestaltet werden. Der Einsatz ist nur dann gerechtfertigt, wenn Werte wie Nachhaltigkeit Transparenz, Fairness und Effizienz nicht verletzt, sondern unterstützt werden.

(Luciano Floridi - Informationsethik) Argumentiert für die moralische Relevanz von nicht menschlichen Informationssystemen. Es erlaubt uns die Abwägung zwischen verschiedenen Informationssystemen, wo ein energieeffizienteres, technisch einfacheres System aus sozialer und ökologischer Sicht gegen den Einsatz von einem ressourcenintensiven, komplexen System spricht \cite{floridiEthicsInformation2015}.

\section{Technikfolgenabschätzung}

\section{Technik-soziologische Perspektiven}

\section{Kriterien für den verantwortlichen KI-Einsatz}

Aus den drei Bereichen lassen sich folgende Kriterien ableiten.

\begin{itemize}
    \item \textbf{Notwendigkeit}: Ist der KI-Einsatz überhaupt erforderlich?
    \item \textbf{Effizienz}: Ist ein einfacheres, energieärmeres System ausreichend?
    \item \textbf{Transparenz}: Kann die Funktionsweise nachvollzogen werden (vom Benutzer und vom Entwickler)?
    \item \textbf{Nachhaltigkeit}: Wie groß sind die ökologischen und sozialen Kosten?
    \item \textbf{Verhältnismäßigkeit}: Steht der Nutzen der KI im Verhältnis zu Aufwand und Folgen?
    \item \textbf{Abhängigkeiten}: Entstehen einseitige Bindungen an andere Systeme oder Anbieter?
    \item \textbf{Werteorientierung}: Unterstützt die Technik gesellschaftliche Werte wie Fairness, Partizipation und Nachhaltigkeit?
\end{itemize}

Diese Kriterien bilden die Grundlage zur Bewertung des Praxisprojekts.